\documentclass[a5paper,10pt]{article}

% https://github.com/InDevRus/filippov-solutions

% Нумерация страниц
\pagenumbering{gobble}

% Настройка полей
\usepackage{geometry}
\geometry{left=1cm}
\geometry{right=1cm}
\geometry{top=1cm}
\geometry{bottom=1.5cm}

% Вставка таблицы
\usepackage{array}
\usepackage{tabularx}

% Инструменты выравнивания формул
\usepackage{amsmath}
\usepackage{mathtools}

% Диакритика в математике
\usepackage{amsfonts}

% Вычисления размеров на ходу
\usepackage{calc}

% Удобные записи производных
\usepackage[italic=true]{derivative}

% Настройки сносок
\usepackage{enotez}

% Длинные знаки векторов
\usepackage{anyfontsize}
\usepackage[b]{esvect}

% Вставка картинок
\usepackage{graphicx}

% Вставка красной строки
\usepackage{indentfirst}
\setlength{\parindent}{1em}

% Условия в командах
\usepackage{xifthen}

% Луа-скрипты
\usepackage{luacode}

% Настройка команд отдельно в математическом и текстовом режимах
\usepackage{mathcommand}

% Для повторения LaTeX кода
\usepackage{multido}

% Конфигурация отступов формул
\delimitershortfall-1sp
\usepackage{mleftright}
\mleftright

% Растягивание объектов
\usepackage{relsize}

% Обтекание текстом
\usepackage{wrapfig}
\usepackage{subcaption}
\usepackage{floatflt}

% Поддержка ввода кириллицы
\usepackage[english, russian]{babel}

% Настройка корневой позиции для компилятора
\usepackage{import}
\providecommand*{\ProjectRootPath}{..}

\import{\ProjectRootPath/preambles/macros/}{theorems}

\import{\ProjectRootPath/preambles/fonts}{font_setup}

% Знак градуса
\usepackage{siunitx}

\import{\ProjectRootPath/preambles/macros/}{operators}
\import{\ProjectRootPath/preambles/macros/}{redefinitions}

\import{\ProjectRootPath/preambles/fonts}{letters_setup}
\import{\ProjectRootPath/preambles/fonts/adjustments}{custom_superscripts}

\import{\ProjectRootPath/preambles/custom}{formatting}
\import{\ProjectRootPath/preambles/custom}{frames}
\import{\ProjectRootPath/preambles/custom}{list_setup}

% TODO: перевести команды на современный стиль объявления
\input{../preambles/plotting}

\begin{document}

$\ddot{x} = 1 - \tpow{e}{x}$.
$\tsys{\dot{x}\br{t} = y \\ \dot{y}\br{t} = 1 - \tpow{e}{x}}$;
$\der{y}{x} = \dvfrac {1 - \tpow{e}{x}} {y}$.
$2y \der{y}{x} = 2 - 2\tpow{e}{x}$.
$\pow{y}{2} = 2x - 2\tpow{e}{x} + C$. \\
В качестве уравнения траектории будем использовать
$\pow{y}{2} + 2\br{\tpow{e}{x} - x - 1} = C$, $C > 0$.

Особой точкой будет $M\br{0; 0}$.
$1 - \tpow{e}{x} = 1 - \br{1 + x + \dfrac {1}{2} \pow{x}{2} + o\br{\pow{x}{2}}} 
= -x + O\br{\pow{x}{2}}$.
$\der{y}{x} = \dfrac {-x + O\br{\pow{x}{2}}} {y}$.
$\pow{\lambda}{2} + 1 = 0$.
$\sub{\lambda}{1,2} = \mp i$.
Поскольку траектории обладают симметрией относительно оси $Ox$,
$M\br{0; 0}$ -- <<центр>>.

\begin{floatingfigure}{0.4\textwidth}
    \begin{tikzpicture}
        \pgfplotsset{
            Graph/.style = {
                thin,
                samples = 500,
                restrict y to domain = -5 : 5,
            },
            DirectionGraph/.style = {
                samples = 20,
                quiver = {scale arrows = 0.075},
                restrict y to domain = -5 : 5,
                -stealth,
            },
            DirectionGraph1/.style = {
                DirectionGraph,
                quiver = {
                    u = {dot_x(x, y) / sqrt(dot_x(x, y) ^ 2 + dot_y(x, y) ^ 2)},
                    v = {dot_y(x, y) / sqrt(dot_x(x, y) ^ 2 + dot_y(x, y) ^ 2)}
                }
            },
            DirectionGraph2/.style = {
                DirectionGraph,
                quiver = {
                    u = {-dot_x(x, y) / sqrt(dot_x(x, y) ^ 2 + dot_y(x, y) ^ 2)},
                    v = {-dot_y(x, y) / sqrt(dot_x(x, y) ^ 2 + dot_y(x, y) ^ 2)}
                }
            },
        }
        \begin{axis}[
            trig format plots = rad,
            width = 6.25cm,
            height = 10cm,
            ylabel=$\dot{x}$,
            ymin = -1.4,
            ymax = 1.6,
            xmin = -1.6,
            xmax = 1.2, 
            xtick = {-10, ..., 10},
            ytick = {-10, ..., 10},
            minor tick num = 3,
            declare function = {
                f(\C, \x) = sqrt(abs(\C - 2 * (exp(x) - x - 1)));
                dot_x(\x, \y) = y;
                dot_y(\x, \y) = 1 - exp(x);
            }]
            
            \foreach \s in {-1, 1} {
                \addplot[Graph, domain = -0.483183 - 0.001 : 0.416221 + 0.0025]
                {\s * f(0.2, x)};
                \addplot[Graph, domain = -0.801218 - 0.0025 : 0.632715 + 0.0025]
                {\s * f(0.5, x)};
                \addplot[Graph, domain = -0.483183 - 0.001 : 0.416221 + 0.0025]
                {\s * f(0.2, x)};
                \addplot[Graph, domain = -1.011228 : 0.757114]
                {\s * f(0.75, x)};
                \addplot[Graph, domain = -1.19825 : 0.857678]
                {\s * f(1, x)};
                \addplot[Graph, domain = -1.534419 : 1.018298]
                {\s * f(1.5, x)};
                \addplot[DirectionGraph1, samples = 2, domain = -1.534419 + 0.01 : -1.534419 + 0.075]
                {\s * f(1.5, x)};
                \addplot[DirectionGraph1, samples = 12, domain = -1.534419 + 0.01 : 1.018298 - 0.05]
                {\s * f(1.5, x)};
            }
        \end{axis}
    \end{tikzpicture}
\end{floatingfigure}
Таким образом, ненулевые решения $x\br{t}$ принимают ограниченный повторяющийся набор \linebreak значений и остаются в некоторой окрестности \linebreak нуля. Более точно,
$\tpow{e}{x} - x \le \dfrac {C} {2} + 1$.

\end{document}
